\chapter{Preliminary Results and Discussion}

Table~\ref{tab:3} reports the completed Stages 0--3 under the shared-frame protocol. The values are validation f-mAP percentages at IoU 0.5 on the same 36,717 frames. Stage 0 uses its own baseline candidates; Stages 1--3 reuse the fixed YOLO candidates. The results are project measurements.

\begin{table}[htbp]
\centering
\caption{Completed-stage validation f-mAP (\%).}
\label{tab:3}
\begin{singlespace}
\begin{tabular}{@{}lrrrr@{}}
\toprule
\textbf{Output} & \textbf{Stage 0} & \textbf{Stage 1} & \textbf{Stage 2} & \textbf{Stage 3} \\
\midrule
Agentness & 35.36 & 65.36 & 65.36 & 65.36 \\
Agent & 16.67 & 35.25 & 35.25 & 35.25 \\
Action & 13.23 & 12.86 & 15.92 & 15.92 \\
Location & 12.77 & 13.36 & 19.60 & 19.60 \\
Duplex & 11.32 & 10.04 & 11.13 & 15.52 \\
Triplet & 7.54 & 6.81 & 6.65 & 9.73 \\
\bottomrule
\end{tabular}
\end{singlespace}
\end{table}

The detector substitution raises agent f-mAP from 16.67\% to 35.25\%. However, improved agent detection does not automatically improve compositions. Stage 1 reaches 6.81\% triplet f-mAP, below the Stage 0 value of 7.54\%, while retaining the baseline-detection matching dependency. This comparison does not isolate that dependency’s contribution from changes in candidates and score construction.

Stage 2 removes that matching dependency and improves action and location f-mAP to 15.92\% and 19.60\%. Its free composition columns reach 11.13\% duplex and 6.65\% triplet f-mAP. This separates two observations: the region features support improved primitive classification, but the independently learned composition scores still leave a substantial gap.

\section{Effect of Learned Composition}

Stage 3 increases duplex f-mAP from 11.13\% to 15.52\% and triplet f-mAP from 6.65\% to 9.73\%, while boxes, agentness, agent, action, and location outputs remain unchanged. The changes are 4.39 and 3.08 percentage points for duplex and triplet, respectively. Relative to Stage 0, Stage 3 is higher on all six output groups in this comparison.

The auxiliary composition comparison provides additional context. An MLP using primitive scores alone reaches 13.37\% duplex and 8.93\% triplet f-mAP; adding the 256-dimensional visual feature raises those values to 15.52\% and 9.73\%. Fixed minimum composition gives 10.36\% and 5.52\%, while corrected single-gate product composition gives 10.58\% and 6.05\%. These measurements support learning composition from both primitive evidence and visual features in this setting. They do not establish a universal advantage over every possible structured model.

\section{Rare-Composition Performance and Limitations}

Table~\ref{tab:4} separates the overall triplet result from the training-count-defined tail. The \(z < 0\) group contains 47 classes and the \(z < -0.5\) group contains 28 of those same classes. These memberships remain fixed across stages.

\begin{table}[htbp]
\centering
\caption{Overall and tail triplet AP (\%).}
\label{tab:4}
\begin{singlespace}
\begin{tabular}{@{}lrrrr@{}}
\toprule
\textbf{Triplet group} & \textbf{Stage 0} & \textbf{Stage 1} & \textbf{Stage 2} & \textbf{Stage 3} \\
\midrule
All 86 & 7.54 & 6.81 & 6.65 & 9.73 \\
\(z < 0\) (47) & 4.14 & 3.43 & 2.60 & 4.53 \\
\(z < -0.5\) (28) & 2.88 & 2.19 & 1.47 & 3.31 \\
\bottomrule
\end{tabular}
\end{singlespace}
\end{table}

Stage 3 recovers the tail losses seen in Stages 1 and 2 and exceeds the Stage 0 primary-tail result by 0.39 percentage points.

The controlled Stage 2-to-3 comparison is particularly informative for the rare-event question: primary-tail AP rises from 2.60\% to 4.53\%, a gain of 1.93 percentage points, while deeper-tail AP rises from 1.47\% to 3.31\%, a gain of 1.84 points. Candidate boxes and primitive outputs are fixed in this comparison, so these gains accompany the learned-composition replacement rather than a new detector. Nevertheless, primary-tail AP remains 4.53\%, compared with 9.73\% over all triplets. The deeper-tail value is 3.31\%. This is the central empirical motivation for Stage 4: gains in overall composition recognition leave rare-event performance low.

The difference between overall and tail means is descriptive. It is not a significance test, nor is it direct evidence that language will solve the remaining problem. The proposed study asks whether semantic structure helps under controls that separate it from representation and classifier changes.

The completed sequence demonstrates that localization, feature access, and composition learning have different effects. Stage 1 changes localization and agent outputs but preserves a detection-matching bottleneck. Stage 2 obtains features for every candidate. Stage 3 changes composition predictions while preserving the earlier outputs. This last transition is the strongest controlled comparison for the learned-composition claim.

These are results on a shared validation population used during development. They do not constitute independent test-set confirmation, a tube-level leaderboard submission, or evidence of real-world driving safety. The completed table also does not establish repeated-seed significance. Accordingly, the proposal treats it as preliminary evidence supporting a focused investigation rather than as a final generalization claim.
